\subsection{How reliable is the forecast? A pre-specified calibration test}\label{sec:forecast}
The forecast was the weakest part of the tool in every evaluation above, so we tested directly whether it can be made reliable. Addendum v3.4 fixed, before any forecaster was fitted, five candidate forecasters, a leave-one-family-out design and five criteria (Supplement Sec.~\ref{sec:S-prereg}). The data are as follows: for each of 44 benchmark problems (\SfDevFamilies{} families) and the \SfHeldProblems{} held-out problems (the eight further problems, the cantilever and the finely meshed plate; four of the families are not in the development set), each seed and each starting budget $N_0\in\{64,128,256\}$, the tool selected a model and returned its cross-validated learning curve; the selected model was then retrained on $2N_0$ and $4N_0$ runs and scored on the test inputs. This gives \SfDevCases{} development forecasts and \SfHeldCases{} held-out forecasts, each with a known outcome. We release these learning curves with the code.
\paragraph{The raw forecast is unbiased but rough} On the development problems the tool's power-law forecast has no bias (median error of the actual relative to the forecast \SfZeroBiasDev{} decades) and is within a factor two in \SfZeroDevShare{} of the cases (\SfZeroHeldShare{} on the held-out problems, median error factor \SfZeroHeldMed). The errors are heavy-tailed: over all \SfAllCases{} forecasts the actual error was within a factor \SfQHalf{} of the forecast in half of the cases, within \SfQEighty{} in 80\% and within \SfQNinety{} in 90\%. It also gets worse with the starting budget: the 80\% quantile of the factor is \SfPhQEightySixtyfour{} from 64 runs, \SfPhQEightyOnetwentyeight{} from 128 and \SfPhQEightyTwofiftysix{} from 256. This is consistent with the failures of H2 above: the share within a factor two (about two thirds) is close to, but under, the threshold of 70\% that we had set.
\paragraph{The pre-specified calibrated forecaster did not help} The proposed forecaster F3 corrects the raw forecast by a ridge regression on ten features of the curve (slopes, curvature, non-monotonicity, error level, number of inputs and outputs, ratio to the POD floor) and adds a conformal interval from the residuals of other families. It met \SfMetCount{} of the five criteria on the held-out problems (\cref{tab:sfforecasters}). Coverage of the nominal 80\% interval was \SfKoneCov{} (problem-level interval \SfKoneLo{} to \SfKoneHi), inside the pre-specified range of 70\% to 90\% (K1 \SfKoneVerdict), and when its upper bound was below a target the target was reached in \SfKfiveOneShare{} of \SfKfiveOneN{} cases for 1\% and \SfKfiveTenthShare{} of \SfKfiveTenthN{} for 0.1\% (K5 \SfKfiveOneVerdict). But F3 was less accurate than the raw forecast (median error ratio \SfKtwoRatio, share within a factor two \SfKtwoGain{} points lower: K2 \SfKtwoVerdict), the constant correction F2 did as well or better on the share and the learned features did not narrow the intervals enough (K3 \SfKthreeVerdict), and the median width of the 80\% interval was \SfKfourWidth{} decades (K4 \SfKfourVerdict). Of our predictions, K2 was wrong (we predicted it would hold); K1, K3, K4 and K5 came out as predicted or, for K1, within the even odds we had given. The learned correction fails because the raw forecast has no bias to remove; its problem is the tails, and a ridge model of ten features cannot predict them. The shrunk-slope forecaster F1 selected no shrinkage and equals the raw forecast.
\begin{table}[t]\centering\caption{Forecasters on the development problems (leave-one-family-out) and on the held-out problems: share of forecasts within a factor two of the actual error, coverage of the nominal 80\% interval and its median width in decades. F0 is the tool's forecast, F5 carries the last observed error forward, F2 adds a constant correction, F3 and F4 are the proposed calibrated forecasters (F4 with a scaled interval).}\label{tab:sfforecasters}\small
\begin{tabular}{lrrrr}\toprule Forecaster & Dev. share & Held-out share & Held-out coverage & Width\\\midrule
F0 raw forecast & \SfZeroDevShare & \SfZeroHeldShare & -- & --\\
F5 carry forward & \SfFiveDevShare & \SfFiveHeldShare & -- & --\\
F2 constant correction & \SfTwoDevShare & \SfTwoHeldShare & \SfTwoHeldCov & \SfTwoDevWidth\\
F3 proposed & \SfThreeDevShare & \SfThreeHeldShare & \SfKoneCov & \SfKfourWidth\\
F4 proposed, scaled interval & \SfFourDevShare & \SfFourHeldShare & \SfFourHeldCov & --\\
\bottomrule\end{tabular}\end{table}
\paragraph{An empirical interval for the raw forecast (post hoc)} Because the learned correction does not help, we gave the raw forecast an empirical interval: plus or minus the 80\% (90\%) quantile of the absolute log error, \SfPhFEighty{} (\SfPhFNinety) in factor. This is post hoc; it was designed after the held-out results were known. Its coverage is \SfPhDevCovEighty{} (\SfPhDevCovNinety) in leave-one-family-out on the development problems and \SfPhHeldCovEighty{} (\SfPhHeldCovNinety) on the held-out problems (problem-level interval for the 80\% level \SfPhHeldLo{} to \SfPhHeldHi), at widths of \SfPhWEighty{} and \SfPhWNinety{} decades. A statement ``the target is reached'' made when the upper end of the 80\% interval is below the target was right in \SfDecOnePrec{} (1\%) and \SfDecTenthPrec{} (0.1\%) of the cases in which it was made, but it found only \SfDecOneRec{} and \SfDecTenthRec{} of the true reaches; the simpler rule ``the forecast itself is below the target'' was right in \SfDecOnePprec{} and \SfDecTenthPprec{} of its statements and found \SfDecOnePrecRec{} and \SfDecTenthPrecRec{} of the reaches. Tool version 0.3.0 prints this interval instead of the earlier range between the three-point and four-point fits, which contained the actual error in under 10\% of the held-out cases. In short, the forecast gives the order of magnitude and a calibrated statement of how rough it is; a decision that needs better than a factor of about four needs more runs, not a better forecast.
